\part{L'expérience humaine dans les entreprises libérées}

Jusqu'ici, nous avons parlé du modèle. Il est temps de parler des gens.

Parce qu'une entreprise libérée, vue de l'intérieur, ce n'est pas un schéma d'organisation : c'est un quotidien. Celui du salarié qui découvre l'autonomie — et parfois son vertige. Celui du dirigeant qui doit détenir le pouvoir sans s'en servir, tous les jours, y compris les mauvais. Celui du client, du fournisseur, du territoire, qui traitent avec une entreprise devenue étrange à leurs yeux.

Cette quatrième partie est celle qui porte le mieux le titre du livre : c'est ici qu'on regarde l'envers du décor humain — les épanouissements réels et les souffrances réelles, sans sacrifier les uns aux autres.

Une règle d'honnêteté pour la traverser : les enthousiastes et les déçus ont existé dans les mêmes entreprises, aux mêmes époques. Toute version de l'histoire qui efface l'un des deux camps est une version fausse.
